% =============================================================================
% Kapitel 3: Methodisches Vorgehen
% ARBEITSGERUEST / REVIEW CANDIDATE
% Stand: 03.09.2026
%
% Dieses Dokument enthaelt bewusst NOCH KEINEN ausformulierten Thesis-Text.
% Die Kommentare bilden die Argumentationsstruktur, Kernaussagen, Quellenbasis,
% Projektevidenz, Abgrenzungen sowie offene Quellenluecken ab.
%
% Quellenstatus:
% [REPO-BIB]  Quelle ist im Thesis-Repository und in literatur.bib vorhanden.
%             Der konkrete Claim muss vor finaler Verwendung im Volltext
%             nochmals geprueft werden.
% [REPO-FILE] Quelle ist als Datei im Repository vorhanden, aber noch nicht
%             verlaesslich in literatur.bib eingebunden und/oder claimweise
%             geprueft.
% [ADD]       Empfohlene Quelle ist noch nicht im Repository und muss vor
%             Verwendung beschafft, gespeichert, gelesen und in literatur.bib
%             aufgenommen werden.
% [PROJECT]   Projektevidenz. Belegt den tatsaechlichen Projektverlauf,
%             ersetzt aber keine wissenschaftliche Literatur fuer allgemeine
%             Aussagen.
% [PRACTICE]  Praxisimpuls ohne empirischen Auswertungsanspruch.
%
% Arbeitsregel:
% Erst nach Pruefung der hier zugeordneten Quellen wird der jeweilige Abschnitt
% in Fliesstext ueberfuehrt.
% =============================================================================

\section{Methodisches Vorgehen}
\label{sec:methodik}

% -----------------------------------------------------------------------------
% FUNKTION DES KAPITELS
% -----------------------------------------------------------------------------
% Zweck:
% Dieses Kapitel erklaert, WIE das Architekturkonzept entwickelt, begruendet,
% konkretisiert und untersucht wurde.
%
% Das Kapitel beschreibt NICHT die finale Architektur im Detail.
% Die resultierende Architektur folgt in Kapitel 4.
%
% Der methodische rote Faden:
%
% wissenschaftliche Ausgangsbasis
% -> literaturbasierte Ausgangshypothese
% -> Forschungsfragen und Engineering-Anforderungen
% -> modellbasierte Architekturentwicklung
% -> Gestaltungsprinzipien fuer KI-gestuetzte Informationsverarbeitung
% -> iterative Architekturentscheidungen
% -> prototypische Realisierung als Untersuchungsmittel
% -> Verifikation und Validierung
%
% Kernaussage des gesamten Kapitels:
% Die finale Informationsverarbeitungsarchitektur entstand nicht aus einer
% einmaligen Implementierungsentscheidung. Sie wurde aus Literatur, fachlichen
% Anforderungen und einer fruehen Architekturhypothese abgeleitet, modellbasiert
% konkretisiert und durch prototypische Untersuchungen iterativ weiterentwickelt.
%
% NICHT HIER:
% - vollstaendige Beschreibung der neun finalen Subsysteme
% - detaillierte Implementierungsbeschreibung einzelner Python-Module
% - konkrete Versuchsergebnisse
% - Interpretation der Versuchsergebnisse
% - Reflexion der KI als Coding-Werkzeug
%
% Diese Inhalte gehoeren in Kapitel 4, 5, 7 bzw. 8.


% =============================================================================
\subsection{Forschungsdesign und Vorgehensmodell}
\label{sec:forschungsdesign}
% =============================================================================

% Zweck:
% Einordnung der Arbeit als gestaltungsorientierte, empirisch rueckgekoppelte
% Architekturuntersuchung. Der Untersuchungsgegenstand ist das
% Architekturkonzept. Der Prototyp dient seiner Konkretisierung und Untersuchung.
%
% Zentrale Frage des Abschnitts:
% Welche Forschungslogik verbindet Literatur, Architekturentwicklung,
% prototypische Realisierung und Evaluation?

Die vorliegende Arbeit verfolgt einen gestaltungsorientierten Forschungsansatz.
Gegenstand der Untersuchung ist die Konzeption einer
Informationsverarbeitungsarchitektur zur kontrollierten Überführung heterogener
Engineering-Bestandsinformationen in SysML v2. Die prototypische
Softwareimplementierung stellt dabei kein eigenständiges Produktziel dar, sondern
dient als Proof of Concept zur Konkretisierung und Untersuchung des
Architekturkonzepts.

Ein solches Vorgehen weist Parallelen zu designorientierten Forschungsansätzen
im Bereich von Informationssystemen und Software Engineering auf. Wieringa
beschreibt Design Science als Verbindung zwischen der Gestaltung eines Artefakts
und der Untersuchung seiner Eignung zur Adressierung eines bestehenden Problems
\cite{Wieringa2014}. Diese Einordnung dient in der vorliegenden Arbeit als
methodischer Bezugsrahmen, ohne dass der Projektablauf als unmittelbare Umsetzung
eines fest vorgegebenen Design-Science-Prozessmodells verstanden wird.

Ausgangspunkt der Architekturentwicklung war eine frühe literaturbasierte
Lösungshypothese. Auf Basis des initial untersuchten Forschungsstands wurden
zunächst drei übergeordnete Betrachtungsebenen unterschieden: ein Data Layer für
die Erfassung und Aufbereitung heterogener Informationen, ein Process Layer für
deren gesteuerte Verarbeitung sowie ein Knowledge Layer für semantische
Harmonisierung, Validierung und wissensgestützte Verarbeitung. Parallel wurde
eine erste Prozessfolge aus Feature Ingestion, Sortierung, ontologischer
Anpassung, Mustererkennung, Human-in-the-Loop und Synthese formuliert. Diese
Struktur stellte keine vorweggenommene Zielarchitektur dar, sondern diente als
konzeptioneller Ausgangspunkt für die weitere Untersuchung. Die zugrunde liegende
Literatur umfasste insbesondere Arbeiten zu LLM-gestützter Modellgenerierung,
modellgetriebener Transformation, Metamodellvalidierung, Ontologien und
semantischer Normalisierung sowie agentenbasierter Informationsverarbeitung
\cite{Pan.2025,Cibrian.2025,Li.2022,Cafe.2013,Kapos.2014,Cibrian.2025b,
Meng.2025,Abonyi.2024,Chis.2025}.

Die Ausgangshypothese wurde nicht unmittelbar implementiert, sondern durch die
Forschungsfragen, fachliche Anforderungen und Stakeholder-Bedürfnisse weiter
konkretisiert. Ergänzende Praxisimpulse, unter anderem aus dem fachlichen
Austausch im Umfeld der MESCONF 2026, dienten zur Schärfung einzelner
Fragestellungen, beispielsweise hinsichtlich Human-in-the-Loop, semantischer
Leitplanken und der Grenzen KI-gestützter Modellgenerierung. Solche Impulse
wurden jedoch nicht als eigenständige empirische Evidenz oder als allgemeiner
fachlicher Konsens gewertet. Allgemeine wissenschaftliche Aussagen werden daher
weiterhin durch Literatur getragen.

Die weitere Architekturentwicklung erfolgte iterativ. Einzelne
Architekturannahmen wurden schrittweise konkretisiert, prototypisch umgesetzt
und anhand definierter Tests und Beobachtungen untersucht. Zeigten sich dabei
Probleme, wurden diese nicht ausschließlich als Implementierungsfehler behandelt.
Lag die Ursache auf konzeptioneller Ebene, etwa in unklaren Verantwortungsgrenzen
oder unzureichend definierten Übergängen zwischen Verarbeitungsschritten, wurde
die zugrunde liegende Architekturentscheidung angepasst. Dieses Wechselspiel
zwischen Gestaltung und Untersuchung entspricht ebenfalls einer zentralen
Charakteristik designorientierter Forschung \cite{Wieringa2014}.

Die Forschungsfragen selbst blieben während dieses Prozesses unverändert und
bildeten den konstanten Bezugspunkt für die Weiterentwicklung der Architektur
und deren spätere Bewertung. Untersucht wird damit nicht die langfristige
produktive Einführung eines fertigen Systems, sondern die Frage, inwieweit die
entwickelte Informationsverarbeitungsarchitektur auf konzeptioneller und
prototypischer Ebene geeignet ist, die aus der Problemstellung und den
Forschungsfragen abgeleiteten Anforderungen zu adressieren.


% =============================================================================
\subsection{Literaturrecherche und Ableitung des Forschungsbedarfs}
\label{sec:literaturmethodik}
% =============================================================================

% Zweck:
% Kompakte und reproduzierbare Beschreibung der Literaturmethodik.
% Die detaillierten Suchstrings, Screeningtabellen und Zwischenstaende werden
% in den Anhang verschoben.
%
% WICHTIG:
% Kapitel 2 enthaelt die INHALTLICHE Synthese des State of the Art.
% Kapitel 3.2 beschreibt die METHODE der Recherche.

Die wissenschaftliche Ausgangsbasis der Arbeit wurde durch eine strukturierte
Literaturrecherche erarbeitet. Ziel war es, den Stand der Technik in den für die
Forschungsfragen relevanten Themenfeldern zu erfassen und daraus sowohl den
Forschungsbedarf als auch eine belastbare Grundlage für die anschließende
Architekturkonzeption abzuleiten. Inhaltlich umfasste der Suchraum insbesondere
SysML und SysML v2, modellgetriebene Transformation, generative KI und Large
Language Models im Systems Engineering, den Umgang mit heterogenen
Bestandsinformationen sowie semantische Normalisierung und Ontologien. Die
Vorgehensweise orientierte sich an etablierten Empfehlungen für systematische
Literaturreviews im Software Engineering \cite{KitchenhamCharters2007}. Der
vollständige Literature Search Plan mit Suchraum, Suchparametern und
Auswahlkriterien ist in Anhang~\ref{app:lsp} dokumentiert.

Die Recherche wurde in Scopus, ScienceDirect, IEEE Xplore und der ACM Digital
Library durchgeführt. Zur Strukturierung des Suchraums wurden die
Forschungsfragen zunächst auf Themenfelder übertragen und anschließend in
Suchbegriffscluster für Modellierungssprachen, KI-Technologien,
MBSE beziehungsweise Architecture as Code sowie Brownfield- und
Heterogenitätsprobleme überführt. Die vollständige Definition der verwendeten
Suchcluster und die datenbankspezifischen Suchstrings sind in
Anhang~\ref{app:lsp} aufgeführt. Ergänzend wurde Scopus AI explorativ eingesetzt,
wobei zusätzlich identifizierte Publikationen vor dem Screening mit den
Ergebnissen der regulären Datenbanksuche abgeglichen wurden.

Die Auswahl der Literatur erfolgte in einem zweistufigen Screening. Nach der
Bereinigung von Dubletten wurden die Treffer zunächst anhand von Titel und
Abstract bewertet. Potenziell relevante Publikationen wurden anschließend einer
Volltextprüfung unterzogen. Die Ein- und Ausschlusskriterien waren bereits im
Literature Search Plan definiert und sind ebenfalls in
Anhang~\ref{app:lsp} dokumentiert. Die Nachvollziehbarkeit von Identifikation,
Screening und Auswahl orientiert sich an den Transparenzprinzipien von PRISMA
\cite{Page.2021}. Der initial dokumentierte Suchlauf umfasste 28 Publikationen;
die vollständigen Screeningentscheidungen und deren Zuordnung zu den definierten
Kriterien sind im Literature Search Report in Anhang~\ref{app:lsr}
wiedergegeben.

Die eingeschlossenen Publikationen wurden anschließend entlang der für die
Forschungsfragen relevanten Lösungsbausteine synthetisiert. Dabei zeichneten
sich insbesondere modellgetriebene Transformation, LLM- und agentenbasierte
Modellgenerierung, metamodellgestützte Validierung, semantische
Normalisierung, Human-in-the-Loop sowie Modularität und Interoperabilität als
wiederkehrende Themenfelder ab. Der ursprüngliche Literature Search Report
dokumentiert diese erste Synthese und ist in Anhang~\ref{app:lsr} enthalten.
Die inhaltliche Einordnung und Aktualisierung des Forschungsstands erfolgt in
Kapitel~2; der Literature Search Report dient dagegen der Nachvollziehbarkeit
des ursprünglichen Recherche- und Syntheseprozesses.

Auf Basis dieser Synthese wurden sowohl der in Kapitel~2 dargestellte
Forschungsbedarf als auch die in Abschnitt~\ref{sec:forschungsdesign}
beschriebene frühe Architekturhypothese abgeleitet. Die Literaturrecherche hatte
damit eine doppelte Funktion: Sie diente einerseits der Einordnung des
wissenschaftlichen Stands und andererseits als Ausgangsbasis für die
anschließende Architekturkonzeption.


% =============================================================================
\subsection{Methodik der modellbasierten Architekturentwicklung}
\label{sec:architekturmethodik}
% =============================================================================

% Zweck:
% Erklaeren, wie aus Forschungsbedarf, Stakeholder-Beduerfnissen und
% Architekturhypothese ein konsistentes Architekturmodell abgeleitet wurde.

Die Architekturentwicklung erfolgte modellbasiert und ausgehend von der in
Abschnitt~\ref{sec:forschungsdesign} beschriebenen konzeptionellen
Ausgangshypothese. Als System of Interest wurde die
Informationsverarbeitungsarchitektur des Turing Generators betrachtet. Die
Systemgrenze trennt dabei die eigentliche Verarbeitung von außerhalb des Systems
liegenden Engineering-Informationsquellen und menschlichen Stakeholdern. Für die
Architekturentwicklung wurden fachliche Anforderungen, Systemfunktionen und
logische Verantwortlichkeiten als unterschiedliche Abstraktionsebenen behandelt.
Diese Trennung unterstützt eine nachvollziehbare Architekturbeschreibung und die
explizite Zuordnung von Stakeholder-Anliegen zu den daraus resultierenden
Architekturelementen \cite{ISO42010_2022}.

Zur Strukturierung der Modellierung wurde das RFLP-Prinzip als methodischer
Referenzrahmen verwendet. Ausgehend von fachlichen Anforderungen wurden zunächst
die erforderlichen Systemfunktionen abgeleitet und anschließend logischen
Verantwortlichkeiten beziehungsweise Komponenten zugeordnet. Für die vorliegende
Arbeit wurde die Betrachtung bewusst auf die Ebenen Requirements, Functional und
Logical Architecture begrenzt. Die Physical-Ebene wurde nicht weiter
ausgearbeitet, da sie für die Untersuchung der Informationsverarbeitungsarchitektur
keinen zusätzlichen methodischen Erkenntnisgewinn erwarten ließ und den
Modellierungsumfang unnötig erweitert hätte. Die resultierende RFL-Struktur
diente damit als kompakter Rahmen zur Trennung von fachlichem Bedarf,
Systemfunktion und logischer Verantwortung. Die schrittweise Ableitung von
Anforderungen über Funktionen hin zu logischen Architekturelementen entspricht
dem Grundgedanken RFLP-basierter Architekturentwicklung, wie er auch in der
Literatur beschrieben wird \cite{LiLockettLawson2020}.

Die fachlichen Anforderungen wurden aus den Forschungszielen, den identifizierten
Stakeholder-Bedürfnissen sowie den aus Literatur und früher Architekturhypothese
abgeleiteten Randbedingungen konkretisiert. Dabei wurde darauf geachtet,
Anforderungen so zu formulieren und zu verknüpfen, dass ihre Herkunft und ihre
Beziehung zu nachgelagerten Funktionen und logischen Verantwortlichkeiten
nachvollziehbar bleiben. Traceability wurde deshalb nicht erst nach Abschluss der
Modellierung ergänzt, sondern als Bestandteil der Architekturentwicklung
behandelt. Die Nachvollziehbarkeit zwischen Ursprung, Anforderung und
nachgelagerten Systemelementen entspricht einem grundlegenden Prinzip des
Requirements Engineering \cite{ISO29148}.

Die Ableitung zwischen den Architekturebenen wurde nicht als starre
Einbahnstraße verstanden. Erkenntnisse auf funktionaler oder logischer Ebene
konnten zeigen, dass vorgelagerte Anforderungen zu grob, redundant oder
unvollständig formuliert waren. In solchen Fällen wurden die betroffenen
Elemente überarbeitet und die Konsistenz zwischen den Ebenen erneut geprüft. Die
modellbasierte Architekturentwicklung verlief damit iterativ zwischen
Requirements, Funktionen und logischen Verantwortlichkeiten. Ziel war nicht eine
möglichst frühe Festlegung der späteren Softwarestruktur, sondern eine
nachvollziehbare Zuordnung zwischen Problemstellung, geforderter Fähigkeit und
fachlicher Verantwortung.

SysML v2 wurde als formale Modellierungssprache zur Repräsentation der
entwickelten Architekturelemente und ihrer Beziehungen verwendet
\cite{SysMLv2}. Die grundlegende Einordnung von SysML v2 sowie die Begründung
seiner Verwendung im Kontext des in dieser Arbeit verfolgten
Architecture-as-Code-Ansatzes erfolgen in Kapitel~2. In diesem Abschnitt steht
dagegen ausschließlich die methodische Nutzung der Sprache für die Entwicklung
und Strukturierung des untersuchten Systems im Vordergrund.

Die detaillierte Ausprägung der daraus resultierenden Anforderungen, Funktionen,
logischen Komponenten und ihrer Beziehungen wird in Kapitel~4 dargestellt. Das
dort erneut aufgegriffene RFLP-Konzept ist von der hier beschriebenen
Modellierungsmethodik zu unterscheiden: Während RFLP in diesem Abschnitt als
Strukturierungsrahmen für die Entwicklung des Turing Generators dient, bildet es
innerhalb der später beschriebenen Informationsverarbeitungsarchitektur selbst
einen Gegenstand der Verarbeitung und Modellbildung.


% =============================================================================
\subsection{Gestaltungsprinzipien der KI-gestuetzten Informationsverarbeitung}
\label{sec:ki_gestaltungsprinzipien}
% =============================================================================

% Zweck:
% Die theoretischen und methodischen Prinzipien beschreiben, anhand welcher
% Kriterien die Informationsverarbeitungsarchitektur gestaltet wurde.
%
% WICHTIGE GRENZE:
% Hier steht, WARUM ein Prinzip verwendet wird.
% Kapitel 4 zeigt, WIE dieses Prinzip in der finalen Architektur umgesetzt ist.

Die KI-gestützte Verarbeitung heterogener Engineering-Informationen basiert auf
mehreren grundlegenden Prinzipien. Dazu gehören die Unterscheidung zwischen
probabilistischer Interpretation und regelgebundener Verarbeitung, die Bindung
von Interpretationen an eine nachvollziehbare Informationsgrundlage, der Umgang
mit unterschiedlichen Interpretationsergebnissen sowie die gezielte Einbindung
menschlicher Entscheidungsinstanzen.

Generative KI kann insbesondere bei semantisch offenen Aufgaben eingesetzt
werden, bei denen Informationen interpretiert, eingeordnet oder in einen neuen
fachlichen Zusammenhang überführt werden müssen. Gleichzeitig erfordern formale
Modellierungsprozesse zusätzliche Prüfmechanismen, da sprachlich plausible
Ergebnisse nicht automatisch strukturell oder syntaktisch valide Modelle
darstellen \cite{Cibrian.2025b,Kapos.2014}. Für KI-gestützte
Engineering-Prozesse ist deshalb eine Trennung zwischen generativer
Interpretation, formaler Validierung und fachlicher Entscheidung von besonderer
Bedeutung.

Human-in-the-Loop-Ansätze adressieren die Einbindung menschlichen Wissens,
menschlicher Kontrolle und expliziter Entscheidungen in KI-gestützte
Verarbeitungsprozesse \cite{MosqueiraRey2023,Lazaros2026}. Dabei kann die Rolle
des Menschen je nach System unterschiedlich ausgeprägt sein und von der
Überwachung automatisierter Prozesse bis zur aktiven fachlichen Entscheidung
reichen. Entscheidend ist dabei nicht allein die Anwesenheit eines Menschen im
Prozess, sondern auch, an welchen Stellen eine Intervention möglich ist und auf
welcher Informationsbasis diese Entscheidung erfolgt.

Eine weitere Grundlage bildet die Nachvollziehbarkeit der verwendeten
Informationsbasis. Bei der Verarbeitung heterogener Wissensquellen ist die
Bindung generierter oder interpretierter Ergebnisse an ihre jeweiligen
Ausgangsinformationen wesentlich, um Herkunft und Kontext der Verarbeitung
erhalten zu können \cite{Chis.2025}. Auf dieser Grundlage können auch mehrere
Interpretationsperspektiven miteinander verglichen werden, ohne den Bezug zur
ursprünglichen Information zu verlieren.

Die folgenden Abschnitte erläutern daher die theoretischen Grundlagen von
probabilistischer und deterministischer Verarbeitung, Source Grounding,
Multi-Persona-Interpretation, Varianz und Stabilität sowie Human Authority und
semantischen Leitplanken.

\subsubsection{Trennung probabilistischer Interpretation und deterministischer Verarbeitung}
\label{sec:det_prob}

% KERNAUSSAGEN:
%
% 1. Nicht jeder Verarbeitungsschritt benoetigt ein LLM.
%
% 2. Semantisch offene Aufgaben koennen probabilistische Interpretation
%    erfordern.
%
% 3. Aufgaben mit eindeutig definierbaren Regeln, Identitaeten, Contracts
%    oder Syntaxpruefungen sollen soweit moeglich deterministisch umgesetzt
%    werden.
%
% 4. LLM-Ausgaben erhalten nicht automatisch Engineering Authority.
%
% 5. Die Architektur trennt deshalb Generierung bzw. Interpretation von
%    Validierung, Entscheidung und deterministischer Weiterverarbeitung.
%
% LITERATUR:
%
% [REPO-BIB] \cite{Cibrian.2025b}
%       Metamodel-driven Validation von SysML-v2-Modellen.
%       Claimweise pruefen fuer:
%       - getrennte Validierung generierter Inhalte
%       - formale Konsistenzpruefung
%
% [REPO-FILE] 2605.18747v1.pdf
%       Ning et al. (2026), Code as Agent Harness.
%       Relevanz:
%       - Code als kontrollierbares Ausfuehrungs- und Verifikationssubstrat
%       - Trennung von generativen Agentenaktionen und mechanischer Pruefung
%       - Multi-Agent-Koordination und Verification
%       Status: arXiv / Preprint. Nicht als alleinige starke Evidenz verwenden.
%       TODO BibTeX.
%
% [ADD / PRUEFEN] Literatur zu Nichtdeterminismus von LLM-Ausgaben und
%       deterministischen Guardrails, falls eine allgemeine Aussage zur
%       Zuverlaessigkeit gemacht wird.
%
% PROJEKTEVIDENZ:
% [PROJECT] Architekturtrenngrenzen im finalen System
% [PROJECT] ADRs zur deterministischen Assembly, Generation und Validation
%
% VORSICHT:
% Nicht behaupten, deterministische Verarbeitung sei generell "korrekt".
% Sie ist reproduzierbar und regelgebunden, kann aber falsche Regeln
% deterministisch ausfuehren.


Generative KI eignet sich insbesondere für Aufgaben, bei denen Informationen
nicht ausschließlich anhand eindeutig formulierter Regeln verarbeitet werden
können. Dazu zählen beispielsweise die Interpretation unstrukturierter Texte,
die semantische Einordnung unterschiedlicher Informationsquellen oder die
Ableitung möglicher Zusammenhänge aus sprachlich beschriebenen Sachverhalten.

Large Language Models erzeugen ihre Ergebnisse probabilistisch auf Grundlage
erlernter Muster und Kontextinformationen. Dies ermöglicht eine flexible
Verarbeitung semantisch offener Eingaben, führt jedoch gleichzeitig dazu, dass
identische oder inhaltlich ähnliche Eingaben nicht zwingend zu vollständig
identischen Ausgaben führen müssen.

Demgegenüber stehen Verarbeitungsschritte, deren zulässiges Verhalten durch
explizite Regeln, formale Modelle oder definierte Schnittstellen beschrieben
werden kann. Hierzu gehören beispielsweise die Prüfung von Identifikatoren und
Datenstrukturen, die Validierung definierter Beziehungen, regelbasierte
Modelltransformationen sowie die syntaktische Prüfung erzeugter Modellartefakte.

Modellgetriebene Transformationsansätze zeigen, dass Transformationen zwischen
formal bekannten Modellrepräsentationen anhand expliziter Regeln beschrieben und
automatisiert ausgeführt werden können
\cite{Kapos.2014,Li.2022}. Im Kontext generierter SysML-v2-Modelle können
zusätzlich Metamodelle und formale Sprachregeln verwendet werden, um strukturelle
beziehungsweise syntaktische Eigenschaften der erzeugten Artefakte unabhängig
vom Generierungsprozess zu prüfen
\cite{Cibrian.2025b,SysMLv2}.

Aus methodischer Sicht ist deshalb zwischen \emph{Interpretation} und
\emph{Verarbeitung} zu unterscheiden. Eine probabilistische Interpretation
beantwortet die Frage, welche fachliche Bedeutung aus einer gegebenen
Informationsgrundlage abgeleitet werden könnte.

Eine deterministische Verarbeitung setzt dagegen voraus, dass die Regeln für
den nächsten Verarbeitungsschritt hinreichend eindeutig definiert sind. Ein
generiertes Ergebnis wird dadurch nicht automatisch zu einer autorisierten oder
validen Engineering-Information.

Semantische Interpretation, formale Transformation, technische Validierung und
fachliche Freigabe stellen daher unterschiedliche Aufgaben mit unterschiedlichen
Entscheidungskriterien dar.

Die Trennung dieser Aufgaben erhöht insbesondere die Nachvollziehbarkeit des
Verarbeitungsprozesses. Wird ein Large Language Model auch für Schritte
eingesetzt, die durch explizite Regeln beschrieben werden könnten, vermischen
sich semantische Interpretation und technische Verarbeitung.

Dadurch wird schwerer erkennbar, ob eine Änderung im Ergebnis aus einer anderen
fachlichen Interpretation oder aus einer Variation innerhalb eines eigentlich
regelgebundenen Verarbeitungsschritts entstanden ist. Die Trennung
probabilistischer und deterministischer Verarbeitung ermöglicht dagegen,
unterschiedliche Fehler- und Unsicherheitsquellen getrennt zu betrachten.

Dabei ist deterministische Verarbeitung nicht mit fachlicher Korrektheit
gleichzusetzen. Ein deterministischer Prozess kann eine fehlerhafte Regel ebenso
reproduzierbar ausführen wie eine korrekte. Der wesentliche Unterschied liegt in
der Reproduzierbarkeit und Prüfbarkeit des Verhaltens.

Während probabilistische Interpretation einen möglichen fachlichen
Bedeutungsraum erschließt, kann regelgebundene Verarbeitung anhand definierter
Bedingungen überprüft und bei identischen Eingaben reproduziert werden.

Für KI-gestützte Engineering-Prozesse ergibt sich daraus ein grundlegendes
Gestaltungsprinzip: Probabilistische Verfahren werden dort eingesetzt, wo
semantische Offenheit oder Interpretationsbedarf besteht.

Sobald ein Verarbeitungsschritt durch explizite Regeln beschrieben werden kann,
sollte dessen Verhalten von der probabilistischen Interpretation getrennt und
möglichst formal beziehungsweise deterministisch ausgeführt werden. Dadurch
bleibt erkennbar, an welchen Stellen im Gesamtprozess Unsicherheit durch
Interpretation entsteht und an welchen Stellen ein Ergebnis auf einer definierten
Verarbeitungsvorschrift beruht.

\subsubsection{Source Grounding und gemeinsame Evidence-Basis}
\label{sec:source_grounding}

% KERNAUSSAGEN:
%
% 1. Semantische Interpretation benoetigt einen eindeutig gebundenen
%    Ausgangsgegenstand.
%
% 2. Deshalb wird zuerst source-grounded Evidence identifiziert und
%    erst danach durch unterschiedliche Personas interpretiert.
%
% 3. Mehrere Personas sollen dieselbe Evidence interpretieren.
%    Andernfalls ist beobachtete Differenz nicht eindeutig auf unterschiedliche
%    Interpretation zurueckzufuehren.
%
% 4. Source Provenance wird nicht nur als nachtraeglicher Audit-Trail
%    behandelt. Sie beeinflusst, welche Informationen in welchem Projekt-
%    und Verarbeitungskontext zulaessig sind.
%
% LITERATUR:
%
% [REPO-BIB] \cite{Chis.2025}
%       Heterogeneous Knowledge Bases und Retrieval-Augmented Generation.
%       Claimweise pruefen fuer:
%       - grounding in externen Wissensquellen
%       - Umgang mit heterogenen Wissensbasen
%
% [REPO-BIB] \cite{Bressan.2000}
%       Kontextrepraesentation und Reasoning.
%       Nur verwenden, falls der konkrete Claim nach Volltextpruefung passt.
%
% [REPO-FILE] Jake Van Clief's Interpretable Context Methodology_260703_235902.pdf
%       Relevanz:
%       - explizite Context Boundaries
%       - stufenweise Informationsweitergabe
%       Status: Preprint / methodische Sekundaerevidenz.
%       TODO BibTeX und genaue Claim-Pruefung.
%
% [ADD] geeignete Primaerliteratur zu Data/Information Provenance im
%       Engineering oder zu provenance-aware AI, falls fuer allgemeine
%       Aussagen ueber Provenance erforderlich.
%
% PROJEKTEVIDENZ:
% [PROJECT] ADR-026
% [PROJECT] ADR-027
% [PROJECT] source-local authority chain
% [PROJECT] project-wide Artifact Identity und Project Fit
%
% NICHT HIER:
% - konkrete Klassen, Manifeste und Fingerprints.
%   Diese folgen in Kapitel 5.

Eine vergleichbare Interpretation mehrerer KI-Instanzen setzt voraus, dass klar
definiert ist, worauf sich diese Interpretation bezieht. Unterschiedliche
Antworten mehrerer Agenten können nur dann sinnvoll als Interpretationsvarianz
betrachtet werden, wenn die zugrunde liegende Informationsbasis kontrolliert
wird. Andernfalls lässt sich nicht unterscheiden, ob eine Abweichung aus einer
anderen fachlichen Interpretation oder lediglich aus unterschiedlichen
Ausgangsinformationen resultiert.

Aus diesem Grund ist zwischen source-grounded Evidence und deren anschließender
Interpretation zu unterscheiden. Zunächst wird festgelegt, welche Information
aus einer konkreten Engineering-Quelle Gegenstand der weiteren Verarbeitung ist.
Erst anschließend können unterschiedliche Perspektiven beziehungsweise Personas
dieselbe Evidence interpretieren. Ansätze zur Nutzung heterogener Wissensbasen
und Retrieval-gestützter generativer Verarbeitung verdeutlichen die Bedeutung
einer expliziten Bindung generierter Ergebnisse an externe Informationsquellen
\cite{Chis.2025}.

Source Provenance beschreibt ergänzend die Herkunft und den
Entstehungszusammenhang einer Information. Provenance-Modelle ermöglichen es,
Beziehungen zwischen Daten, Verarbeitungsschritten und beteiligten Akteuren
nachvollziehbar abzubilden \cite{MoreauMissier2013}. Provenance wird damit nicht
ausschließlich als nachträglicher Audit-Trail betrachtet, sondern bildet bereits
während der Verarbeitung einen relevanten Kontext.

Dies ist insbesondere bei der Verarbeitung mehrerer Engineering-Quellen von
Bedeutung. Interpretationen müssen auf ihren jeweiligen Quellkontext
zurückgeführt werden können, bevor Informationen auf Projektebene miteinander
verglichen oder zusammengeführt werden.

Die Trennung zwischen Evidence und Interpretation schafft zugleich eine
Voraussetzung für die Multi-Persona-Verarbeitung. Mehrere Personas interpretieren
denselben nachvollziehbar gebundenen Evidence-Gegenstand. Beobachtete
Unterschiede können dadurch gezielter als Unterschiede in der Interpretation
betrachtet werden, anstatt sie mit Unterschieden in der zugrunde liegenden
Informationsbasis zu vermischen.

\subsubsection{Multi-Persona-Interpretation, Varianz und Stabilität}
\label{sec:multipersona}

% KERNAUSSAGEN:
%
% 1. Mehrere LLM-Instanzen bzw. Personas koennen unterschiedliche
%    Interpretationspfade oder fachliche Perspektiven erzeugen.
%
% 2. Unterschiedliche Ergebnisse sind nicht automatisch ein Fehler.
%    Varianz kann relevante Unsicherheit oder Interpretationsspielraum
%    sichtbar machen.
%
% 3. Uebereinstimmung ist ebenfalls kein Wahrheitsbeweis.
%    Korrelation oder Konformitaet kann zu gemeinsam falschen Ergebnissen
%    fuehren.
%
% 4. Das Turing-Konzept verwendet deshalb KEIN demokratisches
%    Mehrheitsprinzip als Engineering Authority.
%
% 5. Consensus und Variance sind Informationen fuer die nachfolgende
%    Bewertung.
%
% 6. Wiederholte Aufrufe derselben Perspektive werden als
%    Stabilitaetsevidenz interpretiert und nicht als zusaetzliche Stimmen.
%
% 7. Die konkrete Nutzung von Varianz als Aufmerksamkeits- bzw.
%    Escalation-Signal ist als Designentscheidung der Arbeit zu kennzeichnen,
%    sofern keine Literatur eine allgemeinere Aussage traegt.
%
% LITERATUR:
%
% [ADD] Wang et al. (2023):
%       Self-Consistency Improves Chain of Thought Reasoning in Language Models.
%       ICLR 2023.
%       Relevanz:
%       - diverse reasoning paths
%       - Konsistenz ueber mehrere Samples
%       WICHTIG:
%       Das Verfahren waehlt selbst die konsistenteste Antwort.
%       Turing uebernimmt diesen Entscheidungsmechanismus NICHT ungeprueft.
%
% [ADD] Liang et al. (2024):
%       Encouraging Divergent Thinking in Large Language Models through
%       Multi-Agent Debate.
%       EMNLP 2024. DOI: 10.18653/v1/2024.emnlp-main.992
%       Relevanz:
%       - unterschiedliche Agenten koennen Divergenz erzeugen
%       - Selbstreflexion kann in bestehenden Denkpfaden degenerieren
%
% [ADD] Du et al. (2024):
%       Improving Factuality and Reasoning in Language Models through
%       Multiagent Debate.
%       ICML 2024, PMLR 235.
%       Relevanz:
%       - mehrere Modellinstanzen fuer gegenseitige Pruefung
%       - empirische Verbesserungen auf untersuchten Benchmarks
%
% [ADD] Cui et al. (2026):
%       Free-MAD: Consensus-Free Multi-Agent Debate.
%       Findings of ACL 2026.
%       DOI: 10.18653/v1/2026.findings-acl.1600
%       Relevanz:
%       - Grenzen von Konsens und Majority Voting
%       - Konformitaet und Fehlerfortpflanzung
%       - sehr passend zur Turing-Entscheidung "variance, not votes"
%
% [ADD] Kostka & Chudziak (2026):
%       Controlling Uncertainty and Hallucination Risk in Multi-Agent
%       Fact Verification.
%       UAI 2026, PMLR 337.
%       Relevanz:
%       - Konsens kann bei korrelierten Fehlern falsche Sicherheit erzeugen
%       - steigende Disagreement/Deviation wird als Unsicherheitssignal genutzt
%       - besonders passend fuer Varianz als Attention Signal
%
% [REPO-FILE] 2605.14163v1.pdf
%       Agentic Systems as Boosting Weak Reasoning Models.
%       Moegliche Zusatzquelle fuer Committee-/Verifier-Logik.
%       Status: Preprint. TODO BibTeX und Claim-Pruefung.
%
% PROJEKTEVIDENZ:
% [PROJECT] ADR-026
% [PROJECT] ADR-027
% [PROJECT] Multi-Persona-Manifeste und Stability/Variance-Verarbeitung
%
% VORSICHT:
% - Nicht behaupten, Multi-Persona liefere "die Wahrheit".
% - Nicht behaupten, Mehrheit sei automatisch besser.
% - Literaturergebnisse auf Benchmarks nicht ohne Begrenzung auf
%   Engineering-Modellierung uebertragen.
Die Verwendung mehrerer KI-Perspektiven verfolgt nicht das Ziel, durch eine
Mehrheitsentscheidung automatisch eine vermeintlich richtige Antwort zu
bestimmen. Vielmehr können unterschiedliche Instanzen beziehungsweise Personas
denselben Evidence-Gegenstand aus verschiedenen fachlichen Perspektiven
interpretieren. Multi-Agent-Ansätze zeigen, dass die Kombination mehrerer
Modellinstanzen zu unterschiedlichen Argumentations- und Lösungswegen führen kann
\cite{Du2024}. Daraus lässt sich jedoch nicht ableiten, dass Konsens oder Mehrheit
automatisch fachliche Korrektheit bedeuten.

Für die weitere Betrachtung ist daher die zwischen mehreren Interpretationen
beobachtete Abweichung selbst von Interesse. Unter \emph{Varianz} wird in dieser
Arbeit die fachlich relevante Abweichung zwischen mehreren, auf derselben
Evidence basierenden Interpretationsergebnissen verstanden.

Eine hohe Übereinstimmung zeigt zunächst nur, dass unterschiedliche Perspektiven
zu ähnlichen Ergebnissen gelangen. Sie stellt keinen Wahrheitsbeweis dar. Ebenso
bedeutet eine Abweichung nicht automatisch, dass eine der Interpretationen falsch
ist. Sie kann vielmehr auf Mehrdeutigkeit, unterschiedliche fachliche
Perspektiven oder eine unzureichende Informationsgrundlage hinweisen.

Von dieser Varianz zwischen unterschiedlichen Personas ist die
\emph{Stabilität} einer einzelnen Perspektive zu unterscheiden. Wiederholte
Ausführungen derselben Persona können aufgrund des probabilistischen Verhaltens
des zugrunde liegenden Modells unterschiedliche Ergebnisse erzeugen. Die
Übereinstimmung beziehungsweise Abweichung wiederholter Ausführungen liefert
damit eine Aussage über die Stabilität dieser Interpretation.

Varianz und Stabilität beschreiben somit zwei unterschiedliche Eigenschaften.
Varianz betrachtet Unterschiede zwischen verschiedenen Interpretations-
perspektiven, während Stabilität die Reproduzierbarkeit einer einzelnen
Perspektive über mehrere Ausführungen beschreibt. Wiederholte Ausführungen
derselben Persona sind deshalb nicht als zusätzliche unabhängige Stimmen zu
verstehen.

Beide Größen können als zusätzliche Bewertungsinformation verwendet werden.
Insbesondere kann eine erhöhte Varianz darauf hinweisen, dass ein
Informationsgegenstand nicht eindeutig interpretiert wird und deshalb eine
weitergehende fachliche Bewertung erforderlich sein könnte. Dieser Zusammenhang
ist jedoch nicht als allgemeiner Nachweis fachlicher Unsicherheit zu verstehen.
Varianz und Stabilität liefern Hinweise über das Verhalten der
Interpretationsprozesse, ersetzen jedoch keine fachliche Beurteilung.

\subsubsection{Human Authority und semantische Leitplanken}
\label{sec:human_authority}

% KERNAUSSAGEN:
%
% 1. Human in the Loop wird in dieser Arbeit nicht als beliebige
%    Benutzerinteraktion verstanden.
%
% 2. Der Mensch erhaelt Engineering Authority an definierten
%    Entscheidungs- und Promotionsgrenzen.
%
% 3. Ziel ist nicht die manuelle Kontrolle jedes Verarbeitungsschritts.
%    Menschliche Aufmerksamkeit soll an Stellen eingesetzt werden, an denen
%    fachliche Mehrdeutigkeit, Konflikte oder Authority-Uebergaenge bestehen.
%
% 4. Fuer wirksame Entscheidungen muss der Review-Kontext ausreichend
%    Informationen enthalten, beispielsweise:
%    - zugrunde liegende Evidence
%    - Provenance
%    - Alternativen
%    - Varianz bzw. Unsicherheit
%    - bisherige Entscheidungen
%
% 5. Automatisierte technische Validierung und menschliche fachliche
%    Freigabe sind unterschiedliche Funktionen.
%
% 6. Human Approval darf durch einen bestandenen technischen Check nicht
%    stillschweigend ersetzt werden.
%
% LITERATUR:
%
% [ADD] Lazaros, K.; Vrahatis, A. G.; Kotsiantis, S. (2026):
%       Human-in-the-Loop Artificial Intelligence:
%       A Systematic Review of Concepts, Methods, and Applications.
%       Entropy 28(4), 377.
%       DOI: 10.3390/e28040377
%       Relevanz:
%       - Human Oversight, Feedback und Decision-Making an unterschiedlichen
%         Stellen einer AI Pipeline
%       - Taxonomie von HITL-Ansaetzen
%       - Herausforderungen von Oversight
%
% [ADD] Zhu, L.; Lu, Q.; Ding, M. et al. (2026):
%       Designing meaningful human oversight in AI.
%       AI and Ethics 6, 286.
%       DOI: 10.1007/s43681-026-01147-7
%       Relevanz:
%       - Unterschied zwischen operativer AI Agency und evaluativer
%         menschlicher Agency
%       - Risiko von Rubber-Stamp Oversight
%       - explizite Handover Points, Provenance, Escalation und Override
%       - sehr hohe konzeptionelle Naehe zu Human Authority im Turing Generator
%
% [REPO-BIB] \cite{Cibrian.2025}
%       Claimweise pruefen, in welchem Umfang Human Review / Expert Validation
%       im agentenbasierten SysML-v2-Ansatz enthalten ist.
%
% [PRACTICE] MESCONF 2026
%       Praxisimpuls zu Human Control und Guardrails.
%       Nicht als wissenschaftlicher Wirksamkeitsnachweis verwenden.
%
% PROJEKTEVIDENZ:
% [PROJECT] TF2
% [PROJECT] SR_4 / fruehes Expert Validation Requirement
% [PROJECT] Human Review Gates und Decision Binding
% [PROJECT] ADR-023 final model review/publication
%
% BEGRIFF:
% "Effective Intervention" kann als Zielbegriff aus der Forschungsfrage und
% dem Projekt verwendet werden. Fuer allgemeine Aussagen ueber wirksame
% Human Oversight sollten Lazaros und Zhu herangezogen werden.

% KERNAUSSAGEN:
%
% 1. Ontologien und explizite semantische Modelle koennen Begriffe,
%    Relationen und zulaessige Bedeutungsstrukturen formal bzw. kontrolliert
%    beschreiben.
%
% 2. Fuer heterogene Engineering-Informationen ist dies als moeglicher
%    Mechanismus fuer semantische Harmonisierung und Interoperabilitaet
%    relevant.
%
% 3. Die urspruengliche Forschungsannahme betrachtete Ontologien als zentrale
%    Leitplanke fuer die KI-gestuetzte Transformation.
%
% 4. Die spaetere Architektur untersucht Ontologien nicht isoliert, sondern
%    im Zusammenspiel mit:
%    - Source Provenance
%    - expliziten Contracts
%    - Human Authority
%    - deterministischer Assembly
%    - unabhaengiger Output-Validierung
%
% 5. Eine Ontologie alleine wird NICHT als Garantie gegen Halluzinationen
%    dargestellt.
%
% 6. Wenn die Untersuchung zeigt, dass TF3 urspruenglich zu eng auf
%    Ontologien fokussiert war, wird dies spaeter als Forschungsergebnis
%    diskutiert und nicht durch nachtraegliches Umschreiben von TF3 geloest.
%
% LITERATUR:
%
% [REPO-BIB] \cite{Abonyi.2024}
%       Semantische Technologien, Ontologien und Knowledge Graphs fuer
%       industrielle Datenintegration und Wiederverwendbarkeit.
%
% [REPO-FILE] ISO/IEC 21838-2:2021, Basic Formal Ontology.
%       Fuer Definition und formale Rolle einer Top-Level-Ontologie.
%       TODO BibTeX und genaue Clauses.
%
% [REPO-FILE] Drobnjakovic et al. (2022):
%       The Industrial Ontologies Foundry Core Ontology.
%       Relevanz fuer industrielle semantische Standardisierung.
%
% [REPO-FILE] Orellana & Mandrick (2019):
%       The Ontology of Systems Engineering.
%       Relevanz fuer SE-spezifische Semantik.
%
% [REPO-FILE] Lu et al. (2022):
%       Design Ontology Supporting Model-Based Systems Engineering Formalisms.
%
% [REPO-FILE] Liu et al. (2025):
%       SysDO.
%
% [REPO-FILE] Dong et al. (2026):
%       Insights into ontology-based MBSE.
%
% FUER ALLE REPO-FILE-QUELLEN:
% - BibTeX ergaenzen
% - konkrete Claims im Volltext pruefen
%
% [QUELLENLUECKE]
% Fuer eine allgemeine Aussage "Ontologien reduzieren Halluzinationen"
% ist noch eine spezifische, belastbare Quelle erforderlich.
% Bis dahin nur als Designhypothese bzw. Praxisimpuls formulieren.
%
% PROJEKTEVIDENZ:
% [PROJECT] TF3
% [PROJECT] initialer Knowledge Layer
% [PROJECT] Semantic Governance / Vocabulary / Terminology Decisions
% [PRACTICE] MESCONF 2026

Human-in-the-Loop bezeichnet die gezielte Einbindung menschlichen Wissens,
menschlicher Kontrolle oder expliziter Entscheidungen in einen KI-gestützten
Verarbeitungsprozess. Die konkrete Rolle des Menschen kann dabei unterschiedlich
ausgeprägt sein und von der Überwachung automatisierter Schritte bis zur
fachlichen Bewertung oder Entscheidung reichen
\cite{MosqueiraRey2023,Lazaros2026}.

Für Engineering-Anwendungen ist dabei insbesondere relevant, dass menschliche
Beteiligung nicht auf eine rein formale Bestätigung automatisierter Ergebnisse
reduziert wird. Eine wirksame Intervention setzt voraus, dass die für eine
Entscheidung relevanten Informationen verfügbar sind und eine tatsächliche
Möglichkeit zur Korrektur, Zurückweisung oder Freigabe besteht. Hierzu gehören
beispielsweise die zugrunde liegende Evidence, deren Provenance, alternative
Interpretationen sowie Hinweise auf Varianz oder Stabilität.

Der Begriff \emph{Human Authority} bezeichnet in dieser Arbeit die fachliche
Entscheidungsbefugnis des Menschen an solchen Stellen, an denen ein
KI-generiertes oder KI-interpretiertes Ergebnis in eine autorisierte
Engineering-Information überführt werden soll. Human Authority ist damit von
bloßer menschlicher Beteiligung zu unterscheiden. Entscheidend ist nicht allein,
dass ein Mensch Bestandteil des Prozesses ist, sondern dass ihm an definierten
Entscheidungsgrenzen eine tatsächliche fachliche Entscheidungsfunktion zukommt.

Semantische und ontologische Leitplanken bilden hierzu einen ergänzenden
Mechanismus. Ontologien und explizite semantische Modelle können Begriffe,
Relationen und zulässige Bedeutungsstrukturen beschreiben und dadurch die
Harmonisierung heterogener Informationen unterstützen \cite{Abonyi.2024}.
Insbesondere bei unterschiedlich strukturierten oder unterschiedlich benannten
Engineering-Informationen können solche Leitplanken dazu beitragen, gemeinsame
semantische Bezugspunkte herzustellen.

Die Einhaltung semantischer oder ontologischer Regeln ist jedoch nicht mit
fachlicher Korrektheit gleichzusetzen. Eine formal oder semantisch zulässige
Interpretation kann weiterhin unvollständig oder fachlich unangemessen sein.
Semantische Leitplanken begrenzen damit den möglichen Interpretationsraum und
unterstützen die Erkennung von Inkonsistenzen, ersetzen jedoch weder Source
Grounding noch technische Validierung oder menschliche fachliche Bewertung.

Human Authority und semantische Leitplanken erfüllen somit unterschiedliche,
aber komplementäre Funktionen. Semantische Regeln strukturieren und begrenzen
die automatisierte Verarbeitung, während die fachliche Entscheidungsautorität
an den dafür vorgesehenen Stellen beim Menschen verbleibt.

% =============================================================================
% =============================================================================

\subsection{Iterative Konkretisierung und prototypische Untersuchung}
\label{sec:iterative_konkretisierung}

Die Entwicklung der Informationsverarbeitungsarchitektur erfolgte nicht durch
eine vollständige Vorausplanung bis auf Implementierungsebene. Stattdessen
wurden zunächst die für den Untersuchungsgegenstand wesentlichen
Architekturprinzipien, Systemgrenzen und Verantwortlichkeiten festgelegt. Der
Detaillierungsgrad wurde anschließend jeweils dort erhöht, wo eine weitere
Konkretisierung für die Modellierung, prototypische Umsetzung oder Prüfung
erforderlich wurde.

Dieses Vorgehen entspricht einer progressiven Detaillierung des
Untersuchungsgegenstands. Die Architektur wurde dabei nicht als statisches,
bereits zu Beginn vollständig bestimmbares Ergebnis behandelt, sondern
schrittweise konkretisiert. Designorientierte Forschungsansätze beschreiben
vergleichbare Wechselwirkungen zwischen Gestaltung und Untersuchung, bei denen
ein Artefakt entwickelt, untersucht und auf Grundlage gewonnener Erkenntnisse
weiterentwickelt wird \cite{Wieringa2014}.

Die progressive Konkretisierung ist dabei nicht als bewusste Entwicklung
provisorischer Zwischenlösungen zu verstehen. Architekturentscheidungen und
prototypische Umsetzungen wurden jeweils mit dem Anspruch getroffen, innerhalb
des zu diesem Zeitpunkt verfügbaren Erkenntnisstands tragfähig zu sein. Eine
spätere Anpassung erfolgte dann, wenn zusätzliche Erkenntnisse eine Präzisierung,
Erweiterung oder Revision erforderlich machten.

Auf diese Weise konnten frühzeitig stabile Architekturprinzipien festgelegt
werden, ohne Detailentscheidungen vorwegzunehmen, deren Randbedingungen zu diesem
Zeitpunkt noch nicht ausreichend bekannt waren. Gleichzeitig blieb die
Architektur offen für Erkenntnisse, die erst aus der weiterführenden
Modellierung, der prototypischen Umsetzung oder deren Prüfung hervorgingen.


\subsubsection{Progressive Detaillierung und formative Verifikation}
\label{sec:progressive_detaillierung}

Die zunehmende Detaillierung der Architektur wurde durch formative
Verifikation begleitet. Ziel dieser Prüfungen war nicht ausschließlich der
Nachweis einer fehlerfreien Implementierung. Sie dienten ebenso dazu zu
untersuchen, ob die jeweils konkretisierten Architekturannahmen unter den
betrachteten Bedingungen konsistent und umsetzbar waren.

Zeigte sich bei der Modellierung oder prototypischen Umsetzung, dass eine
Verantwortungsgrenze nicht eindeutig definiert war, ein Informationsübergang
nicht ausreichend spezifiziert wurde oder eine zuvor getroffene Annahme zu
unerwünschten Folgewirkungen führte, konnte die zugrunde liegende
Architekturentscheidung erneut betrachtet werden.

Damit entstand eine Rückkopplung zwischen Architekturkonzeption,
Konkretisierung und Prüfung. Neue Erkenntnisse wurden nicht ausschließlich als
lokale Implementierungsprobleme behandelt, sondern zunächst daraufhin bewertet,
auf welcher Ebene ihre Ursache lag. Konnte die Ursache auf eine unzureichende
Architekturannahme zurückgeführt werden, wurde diese auf der entsprechenden
Abstraktionsebene angepasst.

Dieses Wechselspiel zwischen Gestaltung und Untersuchung entspricht der
Grundlogik designorientierter Forschung, bei der die Entwicklung eines
Artefakts und die Untersuchung seines Verhaltens nicht vollständig voneinander
getrennt erfolgen \cite{Wieringa2014}. Die formative Verifikation hatte damit
eine steuernde Funktion innerhalb der Architekturentwicklung und unterstützte
die schrittweise Erhöhung der Detaillierung.


\subsubsection{Prototyp als Untersuchungs- und Evaluationsinstrument}
\label{sec:poc_evaluation}

Die prototypische Realisierung dient in dieser Arbeit als
Untersuchungsinstrument für die entwickelte Informationsverarbeitungsarchitektur.
Der Prototyp stellt damit kein eigenständiges Forschungsziel dar, sondern
operationalisiert ausgewählte Architekturprinzipien, sodass deren Zusammenspiel
unter kontrollierten Bedingungen untersucht werden kann. Diese Rolle eines
Artefakts als Mittel zur Untersuchung einer entworfenen Lösung entspricht
ebenfalls der Logik designorientierter Forschung \cite{Wieringa2014}.

Im Mittelpunkt steht dabei nicht die vollständige Realisierung eines
produktionsreifen Systems. Untersucht wird vielmehr, ob zentrale Eigenschaften
der Architektur innerhalb des gewählten Proof-of-Concept-Umfangs praktisch
umgesetzt und nachvollziehbar geprüft werden können. Dazu zählen insbesondere
die Verarbeitung heterogener Informationsquellen, die Aufrechterhaltung von
Traceability und Provenance, die Trennung unterschiedlicher
Verarbeitungsverantwortlichkeiten, die Einbindung menschlicher
Entscheidungsinstanzen sowie die Überführung akzeptierter Engineering-
Informationen in formale SysML-v2-Artefakte.

Die Bewertung dieser Eigenschaften erfordert unterschiedliche Formen von
Evidenz. Automatisierte Prüfungen können beispielsweise die Einhaltung
definierter Datenstrukturen, Verträge und technischer Invarianten untersuchen.
End-to-End-Prüfungen ermöglichen dagegen die Betrachtung vollständiger
Informationsflüsse über mehrere Verarbeitungsschritte hinweg. Für erzeugte
SysML-v2-Artefakte kann zusätzlich eine von der Generierung getrennte formale
Validierung gegen die verwendete Modellierungsumgebung beziehungsweise
Sprachregeln erfolgen \cite{Cibrian.2025b,SysMLv2}.

Technische Verifikation ist dabei von fachlicher Bewertung zu unterscheiden.
Ein formal valides Modell ist nicht automatisch ein fachlich korrektes oder
geeignetes Engineering-Modell. Die Untersuchung der Architektur berücksichtigt
deshalb zusätzlich menschliche Review- und Entscheidungsprozesse, deren
theoretische Grundlagen in Abschnitt~\ref{sec:human_authority} beschrieben
wurden.

Die Ergebnisse der prototypischen Untersuchung sind auf den realisierten
Proof-of-Concept-Umfang begrenzt. Sie können zeigen, dass die betrachteten
Architekturprinzipien unter den untersuchten Bedingungen operationalisierbar
sind und wie sie miteinander interagieren. Daraus lässt sich jedoch weder eine
allgemeine Wirksamkeit für beliebige Engineering-Domänen noch eine
Produktionsreife des resultierenden Systems ableiten.
